\chapter{CONCLUSION}

The proposed Decentralized Prepaid Smart Power Grid is intended to provide a software-based approach for prepaid electricity management by integrating blockchain technology, machine learning, backend services, database management, and web-based monitoring. The system is designed to address limitations associated with conventional prepaid electricity management by providing improved transparency, automated balance management, consumption analysis, and user monitoring.

The proposed architecture uses a blockchain smart contract to manage prepaid balances and authorization-related operations. The FastAPI backend is responsible for receiving and processing simulated electricity telemetry, managing application data, and communicating with the blockchain and machine-learning components. The Isolation Forest algorithm is proposed for detecting abnormal consumption patterns, while Linear Regression is used for short-term electricity consumption forecasting.

The system follows a hybrid architecture rather than a fully decentralized architecture. Blockchain provides the decentralized component for prepaid balance and authorization management, while the backend, database, machine-learning module, and web interface provide the remaining application services. This separation allows the proposed system to demonstrate the use of decentralized technology without requiring the complete application to operate on the blockchain.

Since physical electricity meters, sensors, relays, and other electrical hardware are outside the scope of the current project, electricity consumption will be represented using simulated telemetry. The project will therefore focus on evaluating the software architecture and the interaction among its major components in a controlled environment.

The proposed system is expected to demonstrate the feasibility of integrating blockchain-based prepaid management with machine-learning-based consumption analysis and web-based monitoring. The developed prototype may provide a foundation for future extensions involving physical smart-metering devices, real electrical load control, broader blockchain deployment, and real-world testing.

Overall, the project aims to demonstrate a practical software architecture for intelligent prepaid electricity management that combines blockchain-based authorization and balance management with automated analysis and monitoring.